\section*{Glossary of Crystallographic Terms}

\begin{description}
\item[Space Group] A mathematical group describing all symmetry operations (rotations, reflections, inversions, and translations) that leave a crystal structure invariant. There are 230 space groups in 3D and 17 wallpaper groups in 2D.

\item[Wallpaper Group] A 2D space group describing the symmetry of patterns that repeat in two dimensions. Also called plane groups. There are 17 wallpaper groups.

\item[Point Group] The set of rotational and reflectional symmetry operations of a crystal, ignoring translations. The point group describes the rotational symmetry of the crystal class.

\item[Symmetry Operation] A transformation (rotation, reflection, inversion, or translation) that maps the crystal onto itself. Each operation is represented by a rotation matrix $\Rmat$ and a fractional translation vector $\boldsymbol{t}$.

\item[Miller Indices] Integer triplet $(h, k, l)$ that labels a reciprocal lattice vector or a family of crystallographic planes. The indices are the reciprocals of the intercepts of the plane with the crystallographic axes.

\item[Star] A set of symmetry-equivalent reciprocal lattice vectors. A star contains all vectors generated by applying the point group operations to a seed vector, plus their negatives.

\item[Multiplicity] The number of vectors in a star. Determined by the order of the point group divided by the order of the stabilizer subgroup of the seed vector.

\item[Reciprocal Lattice] The Fourier transform of the direct (real-space) lattice. Reciprocal lattice vectors $\qvec = h\boldsymbol{b}_1 + k\boldsymbol{b}_2 + l\boldsymbol{b}_3$ label the Fourier modes in the crystal.

\item[Direct Lattice] The real-space lattice defined by basis vectors $\boldsymbol{a}_1, \boldsymbol{a}_2, \boldsymbol{a}_3$. The crystal structure is a periodic arrangement of atoms on this lattice.

\item[Metric Tensor] A symmetric matrix $\Gmat$ encoding the geometry of the reciprocal lattice: $G_{ij} = \boldsymbol{b}_i \cdot \boldsymbol{b}_j$. Used to compute $|\qvec|^2$ for any Miller index triplet.

\item[Phase Constraint] A relationship between Fourier coefficients of symmetry-equivalent vectors: $c_{\Rmat \qvec} = \exp(2\pi i \, \boldsymbol{t} \cdot \qvec) \, c_{\qvec}$. Ensures the field respects space group symmetry.

\item[Fourier Mode] A single term in the Fourier expansion of the field: $A_n \exp(i \qvec_n \cdot \rvec)$. Each mode corresponds to a reciprocal lattice vector.

\item[SAFBBasis] A collection of valid Fourier modes (stars) for a given space group and lattice. Contains all the information needed to generate a symmetry-adapted field.

\item[Field Generation] The process of computing the real-space field $\psi(\rvec)$ from Fourier coefficients using inverse FFT.

\item[Normalization] The process of scaling the field to a standard range, typically $[0, 1]$, using a tanh function.

\item[Square Norm] The average of $\psi^2$ over the unit cell: $\avg{|\psi|^2} = \frac{1}{V} \int_V |\psi|^2 d\rvec$. Measures the overall amplitude of the field.

\item[Centrosymmetric] A crystal structure is centrosymmetric if it contains an inversion center. In terms of the point group, this means the inversion matrix $(-I)$ is a member.

\item[Inversion at Origin] A special case where the inversion operation maps $\rvec \to -\rvec$ (no translation). This simplifies the phase constraints.

\item[Canonical Key] A lexicographic identifier for a family of planes, computed by sorting the Miller indices and taking the absolute value of the first non-zero element.

\item[Closed Star] A star that contains all symmetry-equivalent vectors, including all negatives. Required for real-valued Fourier series.

\item[BFS Coefficient Solver] A breadth-first search algorithm that propagates phase constraints through the star's relationship graph to determine all Fourier coefficients.

\item[FFTW] Fastest Fourier Transform in the West. A widely-used library for computing FFTs. The SAFB library uses FFTW for efficient field generation.

\item[VTK] Visualization Toolkit. A file format for scientific visualization data. The SAFB library exports fields in VTK XML structured point data format.

\item[SCFT] Self-Consistent Field Theory. A mean-field theory for polymer self-assembly. The SAFB library provides symmetry-adapted initial conditions for SCFT simulations.

\item[Block Copolymer] A polymer consisting of two or more different polymer blocks covalently bonded together. Block copolymers self-assemble into periodic nanostructures.

\item[Gyroid] A bicontinuous minimal surface with Ia-3d symmetry. One of the most studied structures in block copolymer science.

\item[BCC] Body-centered cubic. A crystal structure with space group Pm-3m. Block copolymers can form BCC-like density modulations.

\item[FCC] Face-centered cubic. A crystal structure with space group Fm-3m.

\item[Diamond] A crystal structure with space group Pn-3m. Block copolymers can form diamond-like structures.

\item[Wallpaper Group P1] The simplest 2D space group, containing only the identity operation. No symmetry beyond translation.

\item[Wallpaper Group P6] A 2D space group with 6-fold rotational symmetry. Generates hexagonal patterns.

\item[Wallpaper Group P4] A 2D space group with 4-fold rotational symmetry. Generates square patterns.

\end{description}
